\section{Cross-Theme Synthesis and Research Gap}
\label{sec:synthesis}

The literature on speech-translation systems does not show that either Direct
or Cascaded ST is always more robust. Direct systems remove the intermediate
transcript, while Cascaded systems can use separately developed ASR and MT
components and larger task-specific resources \cite{sperber2020endtoend}.
Previous comparisons also show that the better-performing system can change
depending on language direction, training data, segmentation, reference
translations, and evaluation metrics
\cite{bentivogli2021cascade,anastasopoulos2021iwslt,
etchegoyhen2022casestudy}. Therefore, architecture is only one factor in the
comparison and should not be treated as the only explanation for performance
differences.

The accent-ASR literature establishes one important part of the problem. ASR
performance varies across English accent groups, including both seen and
unseen accents \cite{shi2021aesrc,prabhu2023accentcodebooks,swain2026fair}.
Research on ASR error propagation establishes another important link:
recognition errors, including phonetic errors, can affect the translation
quality of Cascaded systems
\cite{le2017disentangling,sperber2019,pida2026}. However, neither group of
studies directly proves that accent affects final speech-translation quality.
WER is also not available as an equivalent internal metric for standard Direct
ST because Direct systems do not produce an intermediate transcript. In
addition, similar levels of ASR error do not always produce the same amount of
translation error.

The closest existing studies reduce the size of the research gap but do not
completely address it. IWSLT 2024 includes accent-level speech-translation
evaluation, but all systems in its accent challenge were Cascaded systems
\cite{ahmad2024iwslt}. Papi et al. compare Direct, Cascaded, and speech-enabled
systems under accent-related conditions, but their broad benchmark does not use
the matched regional-English design required for this research question
\cite{papi2026hearing}. SeamlessM4T robustness studies examine speaker or
vocal-style variation rather than controlled regional accents, and they do not
include a complete matched Whisper-to-MT Cascaded system
\cite{seamlessm4t2023,seamless2023}. Other work reports accent-related
downstream errors, but accent is treated mainly as a qualitative observation
rather than as the main grouping variable \cite{losttranscription2026}.

The evaluation literature also shows that careful dataset control is
necessary. Common Voice provides self-reported accent labels, while CoVoST 2
provides professionally produced translation references. However, dataset
versions, repeated sentences, speaker structure, and reference construction
can affect how results should be interpreted
\cite{ardila2020commonvoice,wang2021covost2}. Prabhu et al. show that
speaker-disjoint and transcript-aware splits can improve accent evaluation,
although their work focuses on ASR rather than speech translation
\cite{prabhu2023accentcodebooks}. Therefore, Direct and Cascaded systems should
be evaluated using the same utterances and references, while speaker structure
and accent balance should be carefully controlled or reported. Differences in
corpus-level metrics should also be supported by paired statistical analysis
rather than interpreted only from raw BLEU or chrF scores
\cite{papineni2002bleu,post2018clarity,koehn2004significance}.

Overall, the literature supports a focused research gap. No study identified
in this survey combines all of the following in one experiment:

\begin{itemize}
  \item controlled regional-English accent groups;
  \item the same utterances and reference translations;
  \item both Direct and Cascaded ST systems;
  \item final translation quality as the main outcome.
\end{itemize}

This leads to the research question: How does regional English accent
affect the translation quality of Direct and Cascaded speech-to-text
translation systems?

The question focuses on comparative effects under matched evaluation
conditions. It does not assume that accent affects one system more than the
other, or that architecture alone explains any observed difference.
